% !TeX spellcheck = hu_HU
%----------------------------------------------------------------------------
\appendix
%----------------------------------------------------------------------------
\chapter*{\fuggelek}\addcontentsline{toc}{chapter}{\fuggelek}
\setcounter{chapter}{\appendixnumber}
%\setcounter{equation}{0} % a fofejezet-szamlalo az angol ABC 6. betuje (F) lesz
\numberwithin{equation}{section}
\numberwithin{figure}{section}
\numberwithin{lstlisting}{section}
%\numberwithin{tabular}{section}

\section{GCC belső reprezentációk}\label{x:sec:gcc-intr}

Az alábbi listázások a \gls{GCC} belső reprezentációit mutatják be.
Elsődleges célja a szakasznak, hogy prezentálja, hogy a természetesnek tűnő tulajdonsága a Clang \gls{AST}-nek, hogy részletes és könnyen elérhető, nem elterjedt a nagy, releváns implementációk között:
még a nyílt forrásúak, mint a \gls{GCC} sem biztosít ilyen lehetőséget.

Az alábbi kimenetek mind \aref{lst:impl-c}.~lista C kódjából készültek.
Figyelemre méltó, hogy az implicit konverzió {\tt int} típusból {\tt unsigned long} típusba nem jelenik meg: a konkrét 1 érték típusa már az elvárt érték lett.

\begin{lstlisting}[numbers=left,caption={GCC GENERIC reprezentáció},language=c]
> gcc -c -fdump-tree-original=stdout ast-demo.c

;; Function demo (null)
;; enabled by -tree-original


{
  return 1;
}
\end{lstlisting}

\begin{lstlisting}[numbers=left,caption={GCC GIMPLE reprezentáció},language=c]
> gcc -c -fdump-tree-gimple=stdout ast-demo.c
long unsigned int demo ()
{
  long unsigned int D.2771;

  D.2771 = 1;
  return D.2771;
}
\end{lstlisting}

\begin{lstlisting}[numbers=left,caption={GCC SSA reprezentáció},language=c]
> gcc -c -fdump-tree-ssa=stdout ast-demo.c

;; Function demo (demo, funcdef_no=0, decl_uid=2768, cgraph_uid=1, symbol_order=0)

long unsigned int demo ()
{
  long unsigned int D.2771;
  long unsigned int _1;

  <bb 2> :
  _1 = 1;

  <bb 3> :
<L0>:
  return _1;

}
\end{lstlisting}

\begin{lstlisting}
CC=clang CXX=clang++ CFLAGS="-ffunction-sections -fdata-sections" CXXFLAGS="-ffunction-sections -fdata-sections" LDFLAGS="-fuse-ld=mold -Wl,--gc-sections" CFLAGS_FOR_TARGET="-ffunction-sections -fdata-sections" CXXFLAGS_FOR_TARGET="-ffunction-sections -fdata-sections" ../gcc-16.2.0/configure --enable-lto --enable-host-pie --program-suffix=-16 --prefix=/usr/local --enable-compressed-debug-sections=all --disable-multilib --enable-threads=posix --enable-__cxa_atexit --enable-plugins --enable-linker-build-id --disable-werror --disable-nls --enable-default-pie --enable-default-ssp --enable-checking=release --with-linker-hash-style=gnu --disable-sjlj-exceptions  --enable-languages=c,c++,lto --disable-bootstra
\end{lstlisting}

\section{A mért blAST program}\label{x:sec:measured}

Az alábbi program listázásban látható, hogy mi a ténylegesen mért blAST program forrása \aref{ch:perf}.~fejezetben.
Alapvetően semmi különleges nincs benne: az elején a rendszeren található Clang könyvtár helyével inicializálja a blAST könyvtárat, majd az egyik méréshez használt fordítási adatbázist beolvassa.
Az adott méréshez a megfelelő került futtatásra.

Maga a lekérdezés egyszerűen egy diagnosztikába megy, és megtalál minden globális élettartamú változót, amelyik nem egy függvény statikus lokálisa és nem egy {\tt \#include} során került feldolgozásra a futás során, hanem ténylegesen az ellenőrzött forrásokban található.

\begin{lstlisting}[numbers=left,caption={A mérések során használt blAST program}]
use <std4.c4a>
use <blast.c4a>

blast_init resource_dir "/usr/lib/llvm/22/lib/clang/22"

#let ast blast_build_from_compile_commands "/home/bodand/src/build/llvm-23.1.2-build/compile_commands.json"
#let ast blast_build_from_compile_commands "/home/bodand/src/project/blast/cmake-build-rel-clang22/compile_commands.json"
let ast blast_build_from_compile_commands "/home/bodand/src/build/systemd-262.1/build/compile_commands.json"

blast ast
	match into diagnostic "global variable defined"
	such variable declaration
	which global
		 , ~static_local
		 , (expanded in main_file)
		 , nil
end_blast
\end{lstlisting}

